\chapter*{References and Sources}
\addcontentsline{toc}{chapter}{References and Sources}

\begin{plainlanguage}
\textbf{In plain terms:} This book draws on a handful of named, attributable sources (Wariboko's ethnography, the Jakobson--Trubetzkoy distinctive-features story as discussed by Grodzinsky) and a wider, unnamed body of literature consulted while developing the institutional and developmental case studies of Part~III and the transfer literature of Chapter~\ref{ch:transfer}. This chapter lists what can actually be verified and cited properly, and is honest about the limits of that verification rather than presenting borrowed authority dressed as a complete scholarly apparatus --- which would be exactly the kind of failure this book's own Chapter~\ref{ch:fidelity} was built to catch.
\end{plainlanguage}

\section*{Directly cited and attributed in the text}

Wariboko, Nimi. ``Counterfoil Choices in the Kalabari Life Cycle.'' \emph{African Studies Quarterly}, Volume 3, Issue 1 (1999). Center for African Studies, University of Florida. ISSN: 2152-2448. The source for the entire Kalabari case study developed in Chapters~\ref{ch:counterfoil} through \ref{ch:recursion}, including the life-cycle stages, the bibife and chieftaincy-installation examples, and Wariboko's own explicit framing of counterfoil choice as testing cultural ``wisdom'' rather than intelligence --- a distinction this book respects throughout rather than overriding.

Grodzinsky, Yosef. ``How Deeply Human is Language? Chomsky, the Brain, and the\ldots'' Lecture, 2026. The source for the distinctive-features case discussed in Section~\ref{sec:corroboration}, including the auditory-cortex electrode finding, the Lunar Society framing referenced in earlier development of this project, and the Fodor--Hinton modularity anecdote. This is cited from a lecture transcript rather than a peer-reviewed publication; the transcript available to this book did not include complete venue or institutional metadata beyond what the lecture itself states, and this citation should be read accordingly --- as attribution to a specific, identifiable lecture, not as a claim to the editorial standing of a published paper.

Jakobson, Roman, and Nikolai Trubetzkoy. Foundational work on distinctive feature theory, Prague School phonology, 1930s. The theoretical claim discussed in Section~\ref{sec:corroboration} --- that phonemes decompose into abstract binary distinctive features --- is associated with this collaboration and period in the history of linguistics, as discussed by Grodzinsky in the lecture cited above. This book has not traced the claim to a single original paper with confirmed publication details, and the attribution should be read as historically accurate at the level of the school and period (the Prague Linguistic Circle's development of distinctive feature theory through the 1930s, associated most closely with Trubetzkoy's subsequently published \emph{Grundz\"uge der Phonologie}, 1939, and Jakobson's later collaborative work) rather than as a precise bibliographic citation to one source.

``Is AI a Threat?'' Panel discussion, the Royal Institute of Philosophy, 26 June 2026. Chaired by Edward Harcourt (Academic Director, Royal Institute of Philosophy, and Director of the Oxford Institute for Ethics in AI), with panellists Federica Lucivero (BRAID Fellow at the Ada Lovelace Institute and Associate Professor in Ethics of Technology at the Ethox Centre, University of Oxford), Jens Munch (technology and AI strategist, entrepreneur, and investor), Priyanka Suresh (Data Scientist in AI Safety, Google DeepMind), and Shannon Vallor (Baillie Gifford Chair in the Ethics of Data and Artificial Intelligence, Edinburgh Futures Institute, and author of \emph{The AI Mirror}). The source for the closure-illusion illustration of Chapter~\ref{ch:room-for-river}, the delegated-competence discussion of Chapter~\ref{ch:two-axes}, and the deflection case of Chapter~\ref{ch:bad-counterfoils}. Cited from a panel transcript rather than a peer-reviewed publication, with the same caveats given for the Grodzinsky lecture above.

\section*{Consulted in developing the case studies of Part~III and Chapter~\ref{ch:transfer}}

The following sources informed the institutional case study of Chapter~\ref{ch:room-for-river}, the developmental case study of Chapter~\ref{ch:development}, and the transfer literature discussed in Chapter~\ref{ch:transfer}. They are not individually name-checked in the prose of those chapters, which draw on them collectively rather than attributing specific claims to specific papers; they are listed here for transparency and so a reader can verify the underlying material independently. Author names are given where confirmed; where a source's byline could not be confirmed with confidence, only the title, venue, and year are listed rather than a guessed attribution.

\subsection*{Room for the River and Dutch flood governance}

``Room for Rivers: Risk Reduction by Enhancing the Flood Conveyance Capacity of The Netherlands' Large Rivers.'' \emph{Geosciences}, MDPI (2018).

``Room for the River: Innovation, or Tradition? The Case of the Noordwaard.'' Book chapter, 2020.

``Dutch Regional Water Authorities: Institutional Change and Continuity versus Environmental Challenges.'' Conference paper, European Society for Environmental History, Vrije Universiteit Amsterdam (2022).

``River Basin Approach and Integrated Water Management: Governance Pitfalls for the Dutch Space-Water-Adjustment Management Principle.'' Journal article (2006).

``Room for the River, No Room for Conflict.'' Research paper (2018).

``Remaking `Nature': The Ecological Turn in Dutch Water Management.'' Research paper.

``Room for the River.'' \emph{Wikipedia}, accessed 2026. Used for general chronology (the 1993/1995 floods, the 2006 Spatial Planning Key Decision, 2007--2015/2018 construction window) cross-checked against the sources above.

\subsection*{Child language development and category formation}

Horst, J.~S., Oakes, L.~M., and Madole, K.~L. ``Time Course of Visual Attention in Infant Categorization of Cats versus Dogs: Evidence for a Head Bias as Revealed through Eye Tracking.'' \emph{Child Development} 76, no.~3 (2005): 614--631.

``Overextensions in Comprehension and Production Revisited: Preferential-Looking in a Study of Dog, Cat, and Cow.'' \emph{Journal of Child Language}, Cambridge University Press.

``A Computational Theory of Child Overextension.'' \emph{Cognition}, Elsevier (2020).

\subsection*{Self-regulation, executive function, and moral disengagement}

Nash, K., Schiller, B., Gianotti, L.~R.~R., Baumgartner, T., and Knoch, D. ``Electrophysiological Indices of Response Inhibition in a Go/NoGo Task Predict Self-Control in a Social Context.'' \emph{PLOS ONE} 8, no.~11 (2013): e79462.

``Self-Control and Cooperation in Childhood as Antecedents of Less Moral Disengagement in Adolescence.'' \emph{Development and Psychopathology}, Cambridge University Press (2021).

``Annual Research Review: On the Relations among Self-Regulation, Self-Control, Executive Functioning, Effortful Control, Cognitive Control, Impulsivity, Risk-Taking, and Inhibition for Developmental Psychopathology.'' Review article.

``Moral Disengagement.'' \emph{Wikipedia}, accessed 2026. Used for the general framing of Bandura's social-cognitive account of moral disengagement; the underlying theoretical framework originates with Albert Bandura's work on social cognitive theory and moral agency.

\section*{A note on what this list does not include}

This book also draws, in Chapter~\ref{ch:creativity}, on the general shape of the existing creativity research literature --- componential and systems-level accounts distinguishing divergent generation, evaluative judgment, domain expertise, persistence, and social uptake. That chapter is explicit that this reflects general, field-level consensus rather than any specific paper, and no individual source is cited there for the same reason none is listed here: attributing a field's general shape to one paper would itself be a fidelity violation of the kind this book has tried consistently to avoid. A reader wanting to pursue that literature further is better served by a current survey of componential creativity research than by a single citation this book is not in a position to select responsibly.
